\section{A specialization criterion for supported classes}
\label{sec:specialization}

Our goal is to detect a sum of coefficient classes pushed forward from
closed points of a curve.  A vanishing sum would be a localization
boundary.  We first use the two Bott maps of Lemma~\ref{lem:bott} to
express that boundary by ordinary tame symbols.  The divisor boundary on
a smooth model then imposes a relation on the special fiber, where a
homomorphism from its Picard group can detect the contradiction.

\begin{proposition}[Picard specialization]
\label{prop:specialization}
Let $L$ be a finite extension of $\Q_5$ containing $\mu_5$, with valuation
ring $R$ and residue field $k$.  Let $\mathcal X\to\Spec R$ be a smooth
projective morphism of relative dimension one with geometrically integral
fibers
\[
 C=\mathcal X_L,\qquad D=\mathcal X_k.
\]
For a closed point $q$ of $C$, let $r(q)$ be its specialization on $D$:
it is the image of the closed point under the unique extension of
$\Spec\kappa(q)\to\mathcal X$ to the valuation ring of $\kappa(q)$.
For a closed point $r$ of $D$, write $[r]\in\Pic(D)$ for the class of the
divisor $(r)$.  Let $m\geq0$ be an integer, and suppose that:
\begin{enumerate}
\item $U$ is the spectrum of an arbitrary localization of the coordinate
ring of a nonempty affine open of $C$, and $p_1,\ldots,p_m$ are distinct
$L$-rational closed points of $U$;
\item $Z$ is a finite set of closed points of $D$, and every closed point
$q$ of $C$ with $r(q)\notin Z$ belongs to $U$;
\item $B$ is an abelian group killed by $5$, and
$\psi:\Pic(D)\to B$ is a homomorphism satisfying $\psi([r])=0$ for every
$r\in Z$.
\end{enumerate}
Put $r_i=r(p_i)\in D(k)$.  Choose $\beta_L\in K_2(L;5)$ with
$\partial_{\mathrm{coef}}\beta_L=[\zeta]$ for a primitive fifth root of
unity $\zeta$, and let $b$ be multiplication by its restrictions, with
$b^2=b\circ b$.  Let $v_L$ be the normalized valuation of $L$.
For $a_i\in L^\times$ and $e_i\in\F_5$, the condition
\begin{equation}\label{eq:specialization-test}
 \sum_{i=1}^m e_i v_L(a_i)\,\psi([r_i])\ne0\quad\text{in }B
\end{equation}
implies
\begin{equation}\label{eq:specialization-conclusion}
 \sum_{i=1}^m (p_i)_*\bigl(e_i b^2[a_i]_5\bigr)
       \ne0\quad\text{in }G_5(U;5).
\end{equation}
Here $[a_i]_5$ is the class of $a_i$ in
$K_1(L;5)=L^\times/L^{\times5}$, and the scalars in
\eqref{eq:specialization-test} act through the natural $\F_5$-module
structure on $B$.
\end{proposition}

\begin{proof}
The sum in \eqref{eq:specialization-test} is zero when $m=0$, so its
nonvanishing implies $m>0$.  Suppose, for a contradiction, that the sum in
\eqref{eq:specialization-conclusion} is zero.  We first turn this assumed
relation into prescribed ordinary tame symbols on the generic curve.

\smallskip\noindent\emph{Localization and Bott untwisting.}
Put $J=L(C)$ and $A_U=\Gamma(U,\mathcal O_U)$.  The ring $A_U$ is a
localization of the coordinate ring of a smooth affine curve.  It is
therefore a Noetherian regular domain of dimension at most one, with
fraction field $J$.  Since one of the $p_i$ is retained, its dimension is
one.  Every finitely generated torsion $A_U$-module is killed by a nonzero
element; the corresponding quotient of $A_U$ has dimension at most zero
and is Noetherian, hence Artinian.  Such a module has finite length, with simple
factors the residue fields of the retained closed points.  This argument
does not require $U$ to be semilocal or the multiplicative set to be finite.

Quillen localization for the torsion category, followed by d\'evissage and
the coefficient convention in Section~\ref{sec:coefficients}, gives the
exact sequence
\begin{equation}\label{eq:curve-localization}
 K_6(J;5)\longrightarrow
 \bigoplus_{q\in U^{(1)}}K_5(\kappa(q);5)
 \longrightarrow G_5(U;5).
\end{equation}
The second arrow is the sum of coherent pushforwards from the points; the
first is the vector of localization residues.  These are ordinary
localization and d\'evissage for a multiplicative set
\cite[\S4, Theorem~3 and Corollary~2; \S5, Theorems~4--5;
\S7, Proposition~3.2 and paragraph~3.4]{Quillen1973Higher}.
The index set of this direct sum can be infinite, but each vector has finite
support.  Every $q\in U^{(1)}$ is a closed point of
$C$, so $\kappa(q)$ is a finite extension of $L$.

Exactness now supplies $\xi\in K_6(J;5)$ whose residue at $p_i$ is
$e_i b^2[a_i]_5$ and whose other residues in $U$ are zero.  The field $J$
has $\operatorname{cd}_5J\leq3$, and each $\kappa(q)$ has
$\operatorname{cd}_5\kappa(q)=2$.  All these fields are perfect and contain
the same primitive root by restriction from $L$.  Thus Lemma~\ref{lem:bott}
applies to the actual multiplication operations used here: it makes
\[
 \begin{aligned}
 b^2 &:K_2(J;5)\longrightarrow K_6(J;5)
     &&\text{surjective},\\
 b^2 &:K_1(\kappa(q);5)\longrightarrow K_5(\kappa(q);5)
     &&\text{injective}.
 \end{aligned}
\]
Choose $\xi_0\in K_2(J;5)$ with $\xi=b^2\xi_0$.  Write
$\partial_q^{(5)}$ for a coefficient localization residue.  The
compatibility in Lemma~\ref{lem:bott-localization} gives
$\partial_q^{(5)}(b^2\xi_0)=b^2\partial_q^{(5)}\xi_0$, with no sign
because $b$ has even degree.  Injectivity on the residue fields then yields
\begin{equation}\label{eq:untwisted-residues}
 \partial_q^{(5)}\xi_0=
 \begin{cases}
 e_i[a_i]_5,&q=p_i,\\
 0,&q\in U^{(1)}\setminus\{p_1,\ldots,p_m\}.
 \end{cases}
\end{equation}

There is still a coefficient obstruction to representing $\xi_0$
integrally.  Its coefficient boundary belongs to
\[
 K_1(J)[5]=\mu_5(J)=\langle\zeta\rangle.
\]
The last equality holds because $L$ already contains all fifth roots of
unity.  Hence there is $\lambda\in\F_5$ with
$\partial_{\mathrm{coef}}\xi_0=\lambda[\zeta]$.  Let $\beta_U$ and
$\beta_J$ be the restrictions of $\beta_L$.  Naturality gives
$\partial_{\mathrm{coef}}\beta_J=[\zeta]$, while all localization residues
of $\beta_J$ on $U$ vanish because it extends as $\beta_U$.  This uses an
extension over the $L$-scheme $U$, not an extension over $R$.  The class
\[
 \eta=\xi_0-\lambda\beta_J\in K_2(J;5)
\]
therefore has zero coefficient boundary and the same residues as $\xi_0$.
The coefficient exact sequence in Section~\ref{sec:coefficients} supplies
an integral representative $z\in K_2(J)$ of $\eta$.  Write
$\partial_q z\in K_1(\kappa(q))=\kappa(q)^\times$ for its ordinary tame
symbol.  Reduction of localization residues commutes with coefficient
reduction, so \eqref{eq:untwisted-residues} becomes
\begin{equation}\label{eq:integral-tame-symbols}
 [\partial_q z]_5=
 \begin{cases}
 e_i[a_i]_5,&q=p_i,\\
 0,&q\in U^{(1)}\setminus\{p_1,\ldots,p_m\}.
 \end{cases}
\end{equation}
In particular, $\partial_q z$ is a fifth power in $\kappa(q)^\times$ at
every retained point other than the $p_i$.  We have obtained an ordinary
$K_2$-class with the required control of its tame symbols.  It remains to
derive the relation that these symbols must satisfy on $D$.

\smallskip\noindent\emph{The horizontal divisor.}
For each closed point $q$ of $C$, let $R_q$ be the valuation ring of the
unique extension of the valuation of $L$ to $\kappa(q)$, let $k_q$ be its
residue field, and let $v_q$ be its normalized valuation.  The extension
$\kappa(q)/L$ is finite, so $R_q$ is a DVR finite over $R$.  Properness
gives the extension
\[
 j_q:\Spec R_q\longrightarrow\mathcal X
\]
used to define $r(q)$.  This map is finite.  Indeed, its graph in
$\Spec R_q\times_R\mathcal X$ is a closed immersion because
$\mathcal X/R$ is separated, and the projection from this fiber product
to $\mathcal X$ is finite, being the base change of $\Spec R_q\to\Spec R$.
Set
\[
 f_q=[k_q:\kappa(r(q))].
\]

Smoothness over the regular DVR $R$ makes $\mathcal X$ a regular surface,
and its function field is $J$.  The codimension filtration for coherent
$K$-theory gives the following row of the coniveau complex:
\begin{equation}\label{eq:surface-complex}
 K_2(J)\longrightarrow
 \bigoplus_{s\in\mathcal X^{(1)}}\kappa(s)^\times
 \longrightarrow
 \bigoplus_{r\in\mathcal X^{(2)}}\Z.
\end{equation}
The assertion needed here is only that the composite is zero: the first
differential $d_1$ of the coniveau spectral sequence satisfies $d_1^2=0$
\cite[\S7, Theorem~5.4 and equation~(5.5)]{Quillen1973Higher}.
There is exactly one vertical point in $\mathcal X^{(1)}$, the generic
point of the integral curve $D$.  Its residue
$\partial_Dz\in k(D)^\times$ contributes the principal divisor
$\operatorname{div}_D(\partial_Dz)$.  The other codimension-one points are
the closed points $q$ of $C$.

We compute their contributions without assuming that a horizontal closure
is normal.  On $\Spec R_q$, the localization boundary
\[
 K_1(\kappa(q))\longrightarrow K_0(k_q)=\Z
\]
is the normalized valuation $v_q$: a unit has boundary zero, and a
uniformizer $\pi_q$ has boundary $[R_q/\pi_qR_q]=[k_q]$
\cite[V.6.1.2 and V.(6.6)]{Weibel2013}.  Functoriality of coherent
localization for the finite map $j_q$ identifies the horizontal boundary
on $\mathcal X$ with the pushforward of this DVR boundary
\cite[\S7, paragraphs~2.7--2.8, Proposition~3.2 and paragraph~3.4]{Quillen1973Higher}.
At the generic point this pushforward is the identity on $\kappa(q)$.
At the closed point it sends $[k_q]$ to
$[k_q:\kappa(r(q))]\,[\kappa(r(q))]$.  Thus the contribution at $r(q)$ is
exactly $f_qv_q(\partial_qz)$.  There is no additional ramification
factor: the boundary already uses the normalized valuation on
$\kappa(q)$, and the closed-point pushforward contributes only the residue
degree $f_q$.

Define the horizontal divisor
\begin{equation}\label{eq:horizontal-divisor}
 H(z)=\sum_{q\in C^{(1)}}f_qv_q(\partial_qz)\,(r(q)).
\end{equation}
The sum is finite because the first arrow of \eqref{eq:surface-complex}
takes values in a direct sum.  All codimension-two points of the proper
surface lie on $D$.  The identity $d_1^2z=0$ therefore says, as an equality
of divisors on $D$,
\[
 H(z)+\operatorname{div}_D(\partial_Dz)=0.
\]
Consequently
\begin{equation}\label{eq:horizontal-pic-zero}
 [H(z)]=0\quad\text{in }\Pic(D).
\end{equation}

\smallskip\noindent\emph{The Picard test.}
Apply $\psi$ to \eqref{eq:horizontal-pic-zero}.  If $r(q)\in Z$, the
corresponding summand vanishes because $\psi([r(q)])=0$.  If
$r(q)\notin Z$, the hypothesis on $U$ retains $q$.  For such a point other
than the $p_i$, \eqref{eq:integral-tame-symbols} makes
$v_q(\partial_qz)$ divisible by five, so its contribution also vanishes in
$B$.  Finally, $\kappa(p_i)=L$, $k_{p_i}=k=\kappa(r_i)$, and hence
$f_{p_i}=1$ and $v_{p_i}=v_L$.  The remaining congruence in
\eqref{eq:integral-tame-symbols} gives
\[
 v_{p_i}(\partial_{p_i}z)\equiv e_i v_L(a_i)\pmod5.
\]
The identity $\psi([H(z)])=0$ is therefore
\[
 \sum_{i=1}^m e_i v_L(a_i)\,\psi([r_i])=0\quad\text{in }B,
\]
contrary to \eqref{eq:specialization-test}.  This proves
\eqref{eq:specialization-conclusion}.
\end{proof}

\begin{remark}[No Gersten exactness assumption]
\label{rem:no-gersten}
The exactness in \eqref{eq:curve-localization} comes from the torsion-category
localization above, valid for the arbitrary localization $U$.
On the mixed-characteristic surface we use only $d_1^2=0$ in
\eqref{eq:surface-complex}; no Gersten injectivity or exactness is assumed.
\end{remark}
